\section{Tails, moments, and wings}
\label{sec:tails}

The endpoint speeds introduced in \eqref{eq:speeds} have an exact
probabilistic meaning.  This section follows the chain
\[
\text{endpoint speed}\longrightarrow
\text{tail scale}\longrightarrow
\text{finite moments}\longrightarrow
\text{asymptotic smile slope}.
\]

\subsection{Endpoint speeds are tail scales}

\begin{proposition}[Exponential tails]
\label{prop:tails}
For every admissible LQD slice there are constants $K_\pm>0$ such that
\begin{equation}
\Prob(X>x)\sim K_+e^{-x/\lambda_+}\quad(x\to+\infty),
\qquad
\Prob(X<x)\sim K_-e^{x/\lambda_-}\quad(x\to-\infty).
\label{eq:exptails}
\end{equation}
Equivalently,
$\Prob(Y>y)\sim K_+y^{-1/\lambda_+}$ and
$\Prob(Y<y)\sim K_-y^{1/\lambda_-}$.
\end{proposition}

\begin{proof}
On the right, \eqref{eq:asymptotes} gives
$x(z)=m+b_++\lambda_+z+O(e^{-z})$.  Its inverse satisfies
\[
z(x)=\frac{x-m-b_+}{\lambda_+}+O(e^{-z(x)}).
\]
Since $1-\Lambda(z)\sim e^{-z}$,
\[
\Prob(X>x)=1-\Lambda(z(x))
\sim e^{(m+b_+)/\lambda_+}e^{-x/\lambda_+}.
\]
The left tail is identical after reversing signs.
\end{proof}

The gross return therefore has power tails, while the log-return has
exponential tails.  The two endpoint coefficients control extrapolation in a
directly readable unit: log-return distance per $e$-fold reduction in tail
probability.

\subsection{The moment strip}

Every expectation may be written either in rank or log-odds:
\begin{equation}
\E[h(X)]=\int_0^1h(Q(u))\,\dd u
=\int_{-\infty}^{\infty}h(x(z))\rho(z)\,\dd z.
\label{eq:rankintegral}
\end{equation}

\begin{proposition}[Moment strip]
\label{prop:strip}
For an admissible LQD slice,
\begin{equation}
\E[e^{rX}]<\infty
\quad\Longleftrightarrow\quad
-\frac1{\lambda_-}<r<\frac1{\lambda_+}.
\label{eq:strip}
\end{equation}
Consequently the last finite positive moment of $Y$ beyond the first and the
last finite inverse moment are
\[
r_+^*=\frac1{\lambda_+}-1,
\qquad
r_-^*=\frac1{\lambda_-}.
\]
\end{proposition}

\begin{proof}
In \eqref{eq:rankintegral}, the right-tail integrand for $h(x)=e^{rx}$ is
asymptotic to $e^{(r\lambda_+-1)z}$; the left-tail integrand is asymptotic to
$e^{(r\lambda_-+1)z}$.  Integrability gives exactly
\eqref{eq:strip}, with divergence at both boundary values.
\end{proof}

The martingale condition is the point $r=1$.  It lies inside the strip exactly
when $\lambda_+<1$, so normalization, the finite-forward wall, and the first
moment are three forms of the same fact.

The asymmetry of the strip is financially important.  A large $\lambda_-$
may eliminate low-order inverse moments of $Y$ without threatening the
forward; negative log-returns contribute at most one to $Y$.  A comparable
right-tail scale would destroy the first positive moment and hence the
forward itself.  For the long NVDA example, $\lambda_->1$, so even
$\E[Y^{-1}]$ diverges, while $\E[Y]=1$ remains exact.

\subsection{Closed Lee slopes}

Lee's moment formula \citep{Lee2004} relates the critical moments to the
asymptotic slopes of total implied variance:
\begin{equation}
\beta_R=\limsup_{k\to+\infty}\frac{w(k)}k=\Psi(r_+^*),
\qquad
\beta_L=\limsup_{k\to-\infty}\frac{w(k)}{|k|}=\Psi(r_-^*),
\label{eq:lee}
\end{equation}
where
\[
\Psi(r)=2-4(\sqrt{r^2+r}-r),\qquad r\ge0.
\]

\begin{proposition}[Speed-to-slope maps]
\label{prop:leeclosed}
The LQD wing slopes are
\begin{equation}
\beta_R=\frac{2\lambda_+}{(1+\sqrt{1-\lambda_+})^2},
\qquad
\beta_L=\frac{2\lambda_-}{(1+\sqrt{1+\lambda_-})^2}.
\label{eq:leeclosed}
\end{equation}
Both maps increase with their endpoint speed.  For small speeds,
$\beta_\pm\sim\lambda_\pm/2$; as $\lambda_+\uparrow1$, $\beta_R\uparrow2$.
For LQD slices the limsups in \eqref{eq:lee} are ordinary limits.
\end{proposition}

\begin{proof}
Substitute $r_+^*=1/\lambda_+-1$ and $r_-^*=1/\lambda_-$ into $\Psi$ and
rationalize.  For example, with $s=\sqrt{1-\lambda_+}$,
\[
\Psi(1/\lambda_+-1)=2-\frac{4s}{1+s}
=\frac{2\lambda_+}{(1+s)^2}.
\]
The left formula is obtained with $t=\sqrt{1+\lambda_-}$.  Monotonicity and
limits follow from the final expressions.  By \cref{prop:tails}, the return
tails are regularly varying; the Benaim--Friz refinement of Lee's formula
therefore gives convergence rather than only a limsup
\citep{BenaimFriz2009}.
\end{proof}

For completeness, the left algebra is
\[
\Psi(1/\lambda_-)
=2-\frac{4(\sqrt{1+\lambda_-}-1)}{\lambda_-}
=\frac{2\lambda_-}{(1+\sqrt{1+\lambda_-})^2}.
\]
The formulas expose two useful scales.  When $\lambda$ is small, the wing
slope is about half the endpoint speed.  On the right, the map reaches Lee's
model-free ceiling $2$ exactly when the first moment reaches its boundary.
The tail coordinate and the no-arbitrage wing bound meet at the same point.

\begin{figure}[htbp]
\centering
\paperfig{\textwidth}{fig_tails.pdf}
\caption{The tail chain.  Left: endpoint speed determines the last finite
moment; the right map reaches the finite-forward wall at $\lambda_+=1$.
Center: Lee's moment-to-slope function.  Right: the closed compositions in
\eqref{eq:leeclosed}.  Markers show the fitted SPY December 2026 and NVDA
December 2027 slices.}
\label{fig:tails}
\end{figure}

The marked market slices occupy only a small part of the admissible range.
For SPY December 2026 the endpoint speeds are
$\lambda_-=\MacSpyLamL$ and $\lambda_+=\MacSpyLamR$, giving slopes
$\beta_L=\MacSpyBetaL$ and $\beta_R=\MacSpyBetaR$.  The long NVDA slice has
$\lambda_-=\MacNvdaLamL$ and $\lambda_+=\MacNvdaLamR$, with slopes
$\MacNvdaBetaL$ and $\MacNvdaBetaR$.  The heavier single-name tails appear
as larger endpoint speeds before the smile is ever plotted.

\subsection{The limit and the traded wing}

The numbers $\beta_\pm$ describe limits, not quoted strikes.  If
$w(k)=\beta|k|+b+o(1)$, then $w(k)/|k|=\beta+b/|k|+o(1/|k|)$ and may approach
its limit from either side.  The sign of the intercept $b$ is not fixed.

\begin{figure}[htbp]
\centering
\paperfig{0.86\textwidth}{fig_lee.pdf}
\caption{Effective variance slopes for the fitted SPY December 2026 slice.
The left wing lies below its Lee limit at the marked deltas, while the right
wing lies above.  The asymptotic slopes are boundary conditions for
extrapolation; they are not substitutes for prices at finite strikes.}
\label{fig:lee}
\end{figure}

No finite option strip identifies an asymptotic decay rate by itself.  In a
fit, $\lambda_\pm$ are selected jointly by the outer quotes, basis order,
regularization, and endpoint coordinates.  The model makes this choice
low-dimensional and interpretable, but it does not turn extrapolation into
observation.

The finite-strike discrepancy is large enough to matter.  On the SPY slice
in \cref{fig:lee}, the effective right slope is
\MacLeeRatioTenDeltaCall\ times its limit at ten delta and
\MacLeeRatioOneDeltaCall\ times at one delta.  On the left, the corresponding
ratios are \MacLeeRatioTenDeltaPut\ and \MacLeeRatioOneDeltaPut, both below
one.  A limiting slope can therefore overstate one traded wing and understate
the other on the same fitted law.
